\section{One packed word, three ordinary-degree channels}
\label{sec:channels}

The incoming selective encodings improve on full Sidon separation by
isolating only monomials that actually occur in the parent system.  Their
homogenized packed square still mixes constant, linear, and quadratic
source terms.  The new observation is that these three ordinary degrees
can be evaluated separately without introducing another unknown.

Let \(r_1,\ldots,r_{22}\) be the raw quadratic residuals of the
31-positive-variable parent in \code{data/parent_quartic31.json}.
The raw positive coordinates are retained, rather than expanding a
positive-to-nonnegative shift into the source coefficients.  In their
specified order, use the weights
\[
\begin{split}
(w_1,\ldots,w_{31})={}&(104,76,79,60,108,8,99,55,36,54,98,11,7,10,44,32,\\
 &106,22,17,1,3,4,37,89,97,107,39,40,86,45,92),
\end{split}
\]
and, in residual order, the target positions
\[
\begin{split}
(q_1,\ldots,q_{22})={}&(1081,1388,1366,763,850,546,1414,1147,1914,1093,821,\\
 &1861,39,1621,238,759,500,1711,1052,283,1343,368).
\end{split}
\]
Put
\[
 W=108,\quad Q_c=1914,\quad K_c=Q_c+2W+1=2131,
 \quad E_c=K_c+W+2=2241,
\]
\[
 H_c=E_c+W+23=2372.
\]
The 31-coordinate order is exactly
\[
\begin{split}
&(J,F,\alpha,z,f,h,i,o,j,s,w,\tau,\eta,\zeta,y,\gamma,\\
&\hspace{13mm}N,E,P,c,a,X,D,C_2,L,F_2,V,g,I,B_{\rm aux},\Delta).
\end{split}
\]

For a source coefficient \(p_{j,\alpha}\), write \(d=|\alpha|\).
Place its reversed term at exponent \(q_j-\alpha\cdot w\) in
\(A_d(T)\).  Its coefficient is \(2p_{j,\alpha}\) for a constant,
a linear monomial, or a square; for a product of two different coordinates
it is \(p_{j,\alpha}\).  Appendix~\ref{app:channels} gives every
coefficient and exponent.  There are 45 quadratic records, 41 linear
records, and eight constant records.

\begin{lemma}[Degree-separated extraction]
For indeterminates \(z_1,\ldots,z_{31}\), define
\[
 Y_z(T)=\sum_i z_iT^{w_i},\qquad
 \mathcal F(T,y)=A_2(T)y^2+A_1(T)y+A_0(T)+T^{Q_c+1}(y^2+1).
\]
Then, for every \(1\le j\le22\),
\[
 [T^{q_j}]\mathcal F(T,Y_z(T))=2r_j(z_1,\ldots,z_{31}).
\]
\end{lemma}
\begin{proof}
The finite certificate checks, separately for each \(d\in\{0,1,2\}\),
that for every supported monomial \(\alpha\) of degree \(d\) and every
monomial \(\delta\) of exactly that degree,
\[
 q_j-\alpha\cdot w+\delta\cdot w=q_k
 \quad\Longleftrightarrow\quad (j,\alpha)=(k,\delta).
\]
All reversed exponents are nonnegative.  This is verified against all
496 pure quadratic monomials, all 31 linear monomials, and the constant
monomial: \(45\cdot496+41\cdot31+8=23599\) interactions.  The
coefficient of a mixed product in \(Y_z^2\) is two, while that of a
square is one.  The stated coefficient factors give exactly \(2r_j\).
All carrier exponents exceed \(Q_c\).  The supplied checker also
reconstructs the 22 identities symbolically with the native ports left
indeterminate, independently of the interaction count.
\end{proof}

Cross-degree collisions are harmless because a term of \(A_1Y_z\)
is never multiplied by a quadratic monomial, and a term of \(A_2Y_z^2\)
is never multiplied by a constant or a linear monomial.  The three
channels are fixed polynomials evaluated at the same base and the same
single packed input.  No extra existential coordinate is introduced.

\subsection{A positive carrier valid before decoding}

Let \(b\ge2\), \(B=a_Sb^2\), and
\[
 \ell=\sum_iB^{w_i},\qquad y=g+\ell.
\]
The shift by \(\ell\) means that packed nonnegative digits \(n_i\)
decode to the positive parent coordinates \(n_i+1\).  By construction,
the reversed exponents are distinct within each channel.  Since each
coefficient has magnitude below \(M\),
\[
 |A_d(B)|<\frac{MB^{Q_c+1}}{B-1}.
\]
Put \(\varepsilon=M/(B-1)<1/4\).  The elementary inequality
\(y^2+|y|+1\le\tfrac32(y^2+1)\) gives, for \emph{every integer}
\(y\),
\begin{equation}
 \tfrac12B^{Q_c+1}(y^2+1)
 <\mathcal F(B,y)<2B^{Q_c+1}(y^2+1).
 \label{eq:carrier-positive}
\end{equation}
In particular, positivity is available before a ratio bridge has recovered
a bound on \(g\).  A carrier proved positive only on the eventual
decoding range would not justify that proof order.

Define the source registers
\begin{align*}
 N_0&=B^{W+1},& Z_c&=B^{K_c+1},& A_T&=\sum_jB^{q_j},\\
 M_c&=N_0-1-(b-1)\ell,& J_c&=2\mathcal F(B,g+\ell)+A_T,\\
 S_c&=g+N_0J_c,&T_c&=M_c+N_0(B-2)A_T.
\end{align*}
For every \(g>0\), \(S_c>0\).  The base-\(B\) digits of \(M_c\)
are \(B-b\) at the 31 weights and \(B-1\) at the remaining positions
through \(W\), so \(0\le M_c<N_0\) and \(T_c\ge0\), without any
power-of-two assumption on \(b\).

If \(0\le g<4bB^W\), then \(g<N_0\), \(\ell<2B^W\), and
\(g+\ell<5bB^W\).  Equation~\eqref{eq:carrier-positive} gives
\[
 0<\mathcal F(B,g+\ell)<52b^2B^{K_c},\qquad
 0<J_c<105b^2B^{K_c}<Z_c.
\]
Also \((B-2)A_T<Z_c\).  Thus \(0\le S_c,T_c<B^{E_c}\).
The stronger bound needed for the ten-integer adapter follows from
\[
 J_c+(B-2)A_T
 =2\mathcal F(B,g+\ell)+(B-1)A_T
 <(104b^2+1)B^{K_c}<\tfrac14Z_c.
\]
Here \(a_S>420\).  Since \(g+M_c<2N_0\) and \(Z_c\ge4\),
\begin{equation}
 S_c+T_c<N_0(2+Z_c/4)<N_0Z_c=B^{E_c}.
 \label{eq:source-sum-range}
\end{equation}

\subsection{Exact zero extraction with signed borrows}

Assume \(b\) is a power of two.  Then \(B\) and \(N_0\) are powers
of two.  In the range just proved,
\[
 g\mathbin{\&}M_c=0
 \quad\Longleftrightarrow\quad
 g=\sum_i n_iB^{w_i},\qquad 0\le n_i<b.
\]
The tuple is unique.  Put \(z_i=n_i+1\), so \(1\le z_i\le b\).
If \(h_s\) is the coefficient of \(T^s\) in
\(\mathcal F(T,Y_z(T))\), the sum \(z=\sum_i z_i\le31b\) gives
\[
 |h_s|\le (M+1)z^2+Mz+M+1<2048Mb^2=B/4.
\]
The polynomial degree is at most \(K_c\).  Write the signed expansion
\[
 J_c=\sum_{s=0}^{K_c}(2h_s+\epsilon_s)B^s,
 \qquad \epsilon_s=\begin{cases}1&s\in\{q_1,\ldots,q_{22}\},\\0&\text{otherwise.}\end{cases}
\]
Ordinary Euclidean digit normalization has incoming carry always zero
or minus one, by induction from the strict coefficient bound.  At a target,
the normalized digit avoids the binary mask \(B-2\) exactly when
\(h_s=0\): a zero coefficient gives digit zero or one; a positive
coefficient gives a digit at least two; a negative coefficient gives a
normalized digit at least two after borrowing.  This argument explicitly
permits an incoming borrow.

Concatenation at the \(N_0\) boundary and the extraction lemma prove
\begin{equation}
 S_c\mathbin{\&}T_c=0
 \quad\Longleftrightarrow\quad
 \left\{
 \begin{array}{l}
 g=\sum_i n_iB^{w_i},\quad0\le n_i<b,\\
 r_j(n_1+1,\ldots,n_{31}+1)=0\quad(1\le j\le22).
 \end{array}\right.
 \label{eq:channel-zero}
\end{equation}
Bitwise AND is used only in the proof, never in the final equations.

For completeness, any positive parent solution gives nonnegative digits
\(n_i=z_i-1\).  They are not all zero: the fifth shifted residual equals
one at the zero vector.  Choose any sufficiently large power of two
\(b>\max_i n_i\).  All weights are positive, so the resulting
\(g\ge B\ge b\), while \(g<bB^W\), and every carry test passes.
These witnesses exist at arbitrarily large admissible powers of two;
the ten-integer construction may restrict them to \(b=4^s\).
